\scanpage{5}
Plus précisément, le foncteur composé
\diagram{\begin{tikzcd}[ampersand replacement=\&]
D_{\unsure{\mathrm{parf}|\mathrm{prof}}}(A) \arrow[r,"\Phi_{\mathcal L}"] \& D_{\unsure{\mathrm{parf}|\mathrm{prof}}}(B) \arrow[r,"\Phi_{{}^s\mathcal L}"] \& D_{\unsure{\mathrm{parf}|\mathrm{prof}}}(A)
\end{tikzcd}}
est \inserted{can.} isomorphe au foncteur \deleted{identique, \illegible{2}}\\
\inserted{$\mathbf Ls_A^*$ \deleted{$\unsure{\mathbf Rs}$} $\simeq\mathbf Rs_{A*}$ suivi de $\tau_A[-n]$,}\\
\deleted{\unsure{respectif}}, \inserted{où $\tau_A\simeq\overset{\unsure{n}}{\Lambda}t_A\simeq\overset{\unsure{n}}{\Lambda}t_B$. Donc} $\Phi_{\mathcal L}$ et $\mathbf Ls_A^*\Phi_{{}^s\mathcal L}$ \inserted{$=\mathbf Rs_{A*}\Phi_{{}^s\mathcal L}$} sont des équivalences de\\
catégories triangulées de $D_{\unsure{\mathrm{parf}|\mathrm{prof}}}(A)$ et $D_{\unsure{\mathrm{parf}|\mathrm{prof}}}(B)$.
\margin{\deleted{\textcircled{\unsure{de catégories}}}}
\margin{inutile de\\
se borner aux\\
complexes parfaits,\\
on peut les prendre\\
p. ex. dans\\
$D^b_{\mathrm{qcoh}}$ \deleted{$\unsure{A}$}}

\uline{Corollaire} \unsure{3|2} $\Phi_{\mathcal L}\cdot\omega_A[n]:$ \deleted{transforme}\\
\inserted{$F\longmapsto\Phi_{\mathcal L}(F\otimes\omega_A[n])\simeq\Phi_{\mathcal L}(F)\otimes\omega_A[n]$ transforme}\\
$\overset{\mathbf L}{\otimes}$ en $\overset{\mathbf L}{*}$ : \deleted{(et \unsure{réciproquement}}\\
\deleted{$\overset{\mathbf L}{*}$ en $\overset{\mathbf L}{\otimes}$)}. $\unsure{D_{\mathcal L}}=$ \deleted{$\unsure{\Phi_{\mathcal L}}$}\\
\inserted{$\deleted{x\mapsto\varphi_{\mathcal L}(\unsure{x})}\quad x\longmapsto(-1)^n\operatorname{cl}\omega_A\,\varphi_{\mathcal L}(x)$} transforme produit ordinaire\\
en \deleted{\unsure{couple}} produit de Convolution.
\[
\Phi_{\mathcal L}(F\unsure{\overset{\mathbf L}{\otimes}|\overset{\mathbf L}{*}}G)\otimes\omega_A[n]
\simeq[\Phi_{\mathcal L}(F)\otimes\omega_A[n]]*[\Phi_{\mathcal L}(G)\otimes\omega_A[n]]
\]
i.e.
\[
\Phi_{\mathcal L}(F\overset{\mathbf L}{\otimes}G)\simeq[\Phi_{\mathcal L}(F)*\Phi_{\mathcal L}(G)]\otimes\omega_A[n].
\]
Prouvons la prop. 2. Il suffit de définir un isom
\[
\Phi_{{}^s\mathcal L}\cdot\Phi_{\mathcal L}\simeq\operatorname{id}_{D_{\unsure{\mathrm{parf}|\mathrm{prof}}}(A)}.
\]
Or le composé des foncteurs\\
est de la forme $F\longmapsto\mathbf R\operatorname{pr}_{2*}(\operatorname{pr}_1^*(F)\overset{\mathbf L}{\otimes}\mathcal M)$, où\\
$\mathcal M\in D_{\unsure{\mathrm{parf}|\mathrm{prof}}}(A)$ est donné par
\[
\mathcal M=\unsure{\mathbf R|\mathbf L}\operatorname{pr}_{13*}\bigl(\operatorname{pr}_{12}^*(\mathcal L)\operatorname{pr}_{23}^*({}^s\mathcal L)\bigr).
\]
\margin{$A\times B\times A$}
Il faut donc calculer ce $\mathbf R\operatorname{pr}_{13*}$, et pour ceci,\\
posons $C=A\times A$, et considérons l'isom\\
$A\times B\times A\simeq\deleted{\unsure{C\times B\times}}C\times B$. Le module inversible\\
$\mathcal N=\operatorname{pr}_{12}^*(\mathcal L)\operatorname{pr}_{23}^*({}^s\mathcal L)$ sur $C\times B$ est \deleted{\unsure{muni de}} birigidifié de façon\\
canonique, et correspond à un homomorphisme\\
$C\xrightarrow{\psi}B^*$, qui n'est autre que l'hom. de $A\times A$ dans $B^*$\\
dont les \inserted{deux} composantes sont égales à $\varphi:A\mathbin{-}B$ (l'hom\\
de $A$ dans $B^*$ défini par $\mathcal L$).
\mspage{4}
